\chapter{Integración Continua, Despliegue y DevOps}
\label{cap:cicd}

Para facilitar la validación de las decisiones de diseño en la IU y de los prototipos funcionales, durante el proyecto se han integrado los conceptos de CI/CD, lo que ha permitido agilizar el despliegue de la aplicación durante todo el proceso de desarrollo. En este sentido, en este capítulo se explica cómo se ha realizado la automatización de la entrega continua de software y la infraestructura para los diferentes entornos de \textit{ClimbIt}, empleando metodologías basadas en la filosofía de DevOps para integrar los nuevos cambios de manera eficaz.

Dada la carga técnica de estos procesos, este capítulo se estructura en tres bloques principales para guiar su lectura:

\begin{itemize}
    \item \textbf{Diseño de la Integración Continua:} Se detalla cómo se evalúa la calidad del código y se automatizan las pruebas previas a la subida de cualquier cambio.
    \item \textbf{Entornos e Infraestructura:} Se describen los espacios de desarrollo, preproducción y producción, abordando la decisión de utilizar un servidor físico propio (\textit{self-hosted}) y la orquestación de la arquitectura mediante contenedores.
    \item \textbf{Empaquetado y subida a Google Play:} Finalmente, se aborda la decisión poco habitual de distribuir la aplicación web ClimbIt a través de la plataforma de Google Play. Se explicará el porqué de esta necesidad operativa para reducir la fricción en la instalación de los usuarios, detallando la tecnología (\textit{Trusted Web Activity}) utilizada para lograrlo.
\end{itemize}
\section{Diseño de la Integración Continua}

La integración continua (CI) es una práctica que permite fusionar nuevo código desarrollado en producción de forma segura y eficiente. En \textit{ClimbIt} se han implementado dos mecanismos para lograr esta integración.

\subsection{Calidad de Código y Validación Pre-commits}

Previo a la subida del código al repositorio remoto en GitHub, se realiza una revisión del código en local en cada \textit{commit} para asegurar que cumple con las normas establecidas en el proyecto. Para ello, se han empleado las siguientes herramientas:

\begin{itemize}
    \item \textbf{Formateo automático con Prettier:} Se utiliza \textbf{Prettier} para que todo el código desarrollado tenga el mismo estilo visual. Se encarga automáticamente de poner puntos y coma al final de las líneas, usar comillas dobles o evitar que las líneas de texto sean demasiado largas.
    
    \item \textbf{Revisión de errores con ESLint:} \textbf{ESLint} es una herramienta que notifica si se están cometiendo errores de lógica o usando malas prácticas al programar. Se ha configurado desactivando alertas innecesarias y avisando si se crean variables que luego no se referencian.
    
    \item \textbf{Automatización de los pre-commits con Husky:} \textbf{Husky} se usa para la configuración de disparadores (\textit{hooks}) de Git que obligan al código a pasar por Prettier y ESLint justo antes de que los desarrolladores realicen un \textit{commit} de los cambios. Si el código falla en la validación de Prettier o ESLint, Husky bloquea el guardado asegurando que solo código de buena calidad se suba al repositorio.
\end{itemize}

\subsection{Workflow de Integración (ClimbIt-CI)}

Una vez que los cambios han pasado las validaciones locales y se suben a GitHub, al crearse una \textit{Pull Request} a una de las ramas principales (\texttt{main} o \texttt{develop}), a través de GitHub Actions se ejecuta el \textit{pipeline} (\textit{ClimbIt-CI}) que realiza los siguientes pasos: 

\begin{itemize}
    \item \textbf{Creación del entorno:} Se crea en un servidor de GitHub un entorno temporal con la base de datos \textit{PostgreSQL} exclusivamente para realizar las pruebas.
    
    \item \textbf{Instalación de dependencias:} Dentro del sistema se ejecuta el comando \texttt{npm ci} para descargar exactamente las mismas librerías que necesita el proyecto.
    
    \item \textbf{Pruebas de Integración:} Prepara la base de datos y ejecuta las pruebas automáticas con la herramienta \textit{Jest} para comprobar que toda la lógica de la aplicación y las conexiones (APIs) funcionan correctamente.
\end{itemize}

\section{Entornos de desarrollo, preproducción y producción}

Para poder desarrollar de forma fluida y segura, se ha dividido el proyecto en tres entornos distintos:

\begin{itemize}
    \item \textbf{Entorno de Desarrollo Local:} Se trata del entorno de desarrollo local en el que los cambios se visualizan al instante (\textit{Hot Reload}). Solo se ejecutan en contenedores de \textit{Docker} la base de datos \textit{PostgreSQL} y su panel de control (\textit{pgAdmin}). El \textit{frontend} y el \textit{backend} se ejecutan localmente con el comando \texttt{npm run dev}.
    
    \item \textbf{Entorno de Preproducción (\texttt{dev.climbit.es}):} Se ejecuta en un \textit{self-hosted runner} a través de GitHub, que se define más adelante. Este entorno tiene la finalidad de validar el funcionamiento de los cambios en un servidor real antes de publicarlos en producción (se ejecuta bajo el proyecto \texttt{climbit\_dev}). Es accesible desde cualquier navegador web\footnote{Entorno de preproducción: https://dev.climbit.es}.
    
    \item \textbf{Entorno de Producción (\texttt{app.climbit.es}):} También se ejecuta en un \textit{self-hosted runner} a través de GitHub y se trata de la versión final, estable y pública de la aplicación, que es la que usan los usuarios finales. Se ejecuta bajo \texttt{climbit\_prod}. Es accesible desde cualquier navegador web\footnote{Entorno de producción: https://app.climbit.es/}.
\end{itemize}

\section{Despliegue de la infraestructura}

Para publicar \textit{ClimbIt}, se tomó la decisión de no usar servidores \textit{Cloud}, ya que realizándolo localmente se obtiene un control total sobre la información además del ahorro económico. Para ello, se preparó un servidor (\textit{self-hosted}) que se puede utilizar mediante GitHub Actions.

\subsection{Orquestación mediante Contenedores}

Tanto el entorno de preproducción como el de desarrollo se organizan usando \textit{Docker Compose}. Esta arquitectura agrupa cuatro contenedores que se conectan entre sí:

\begin{enumerate}
    \item \textbf{Frontend:} Servidor web que entrega la interfaz de usuario.
    \item \textbf{Backend:} API Node.js que gestiona la lógica de negocio.
    \item \textbf{PostgreSQL:} Motor de base de datos relacional.
    \item \textbf{pgAdmin:} Interfaz gráfica para la administración de la base de datos.
\end{enumerate}

\subsection{Hardware y Ejecución (Self-hosted Runner)}

Físicamente, el servidor se encuentra alojado en una \textbf{Raspberry Pi 4} (con 4GB de RAM) usando el sistema operativo \textit{PiOS} basado en \textit{Debian}. Dentro del servidor se ha instalado un \textit{GitHub Self-hosted Runner} como servicio. Se trata de un programa que espera continuamente órdenes de despliegue automático desde GitHub. Cuando recibe la orden, el servidor vuelve a montar la aplicación y le asigna los puertos correspondientes (el puerto 80 para producción y el 81 para preproducción) de forma automática y segura.

\subsection{Exposición Segura y Cifrado}

Para acceder al servidor por internet sin comprometer la red local, se tomó la decisión de utilizar un túnel seguro \textbf{Cloudflare Zero Trust}. Este sistema crea una conexión directa y protegida (\textit{tunnel}) desde los subdominios (\texttt{app.climbit.es} y \texttt{dev.climbit.es}) en los servidores de Cloudflare hasta el servidor donde se encuentra alojado \textit{ClimbIt}.

Toda la información viaja encriptada usando el protocolo \textbf{HTTPS} y Cloudflare se encarga automáticamente de gestionar los certificados de seguridad (\textbf{certificados SSL/TLS}) en sus propios servidores. De esta forma, se garantiza la seguridad de los datos de los usuarios sin tener que lidiar con configuraciones de seguridad complejas.

\section{Empaquetado y Subida a Google Play}

Durante las pruebas realizadas con usuarios finales a lo largo del desarrollo del proyecto, se determinó que existía un subgrupo de escaladores con un perfil tecnológico básico, los cuales se encontraron con una serie de problemas al instalar la aplicación desde navegadores como \textit{Google Chrome} o \textit{Safari}. En base a este suceso, se tomó la decisión de preparar la aplicación para subirla a Google Play, reduciendo de esta manera la fricción de los usuarios al instalarla. 

\subsection{Tecnología Trusted Web Activity (TWA) y Bubblewrap}

Para transformar una PWA (\textit{Progressive Web App}) en un formato que acepte Android (.aab), se utilizó la herramienta \textbf{\textit{Bubblewrap}} desarrollada por Google. Esta herramienta envuelve la aplicación usando la tecnología \textit{Trusted Web Activity} (TWA) y, gracias a esto, la aplicación se asimila a cualquier otra aplicación nativa instalada en el móvil, aunque internamente sigue funcionando conectada al servidor. Tiene la ventaja de que cualquier cambio realizado se actualiza instantáneamente en la aplicación del dispositivo.

\subsection{Configuración de Verificación y Requisitos Legales}

Para poder publicar en Google Play se completaron los siguientes pasos técnicos y legales:
\begin{itemize}
    \item \textbf{Verificación del dominio:} Se configuró en el archivo \texttt{/assetlinks.json}, ubicado en la carpeta \texttt{well-known}, el envío de información para demostrar a Google la propiedad de la aplicación. Esto permite que el dispositivo móvil confíe en la aplicación y la muestre en pantalla completa.
    
    \item \textbf{Política de Privacidad:} Se redactó y publicó una página explicando cómo se tratan los datos, ya que es un requisito para subir aplicaciones a Google Play. Para ello se utilizó una plantilla de la plataforma Iubenda, como se puede observar en la siguiente referencia\footnote{Política de privacidad: https://www.iubenda.com/privacy-policy/86624965}.
    
    \item \textbf{Manifiesto de Aplicación Web (\textit{manifest.json}):} Se configuró este archivo con los iconos, colores y datos de la aplicación para que, al instalarla, se integre visualmente en cualquier dispositivo Android.
    
    \item \textbf{Configuración del perfil de la aplicación:} Además de esto, se tuvo que crear una descripción, un banner, un formulario de eliminación de datos y capturas en dispositivos para la ficha de la tienda.
\end{itemize}

\subsection{Fase de Lanzamiento: Prueba Cerrada}

Actualmente \textit{ClimbIt} se encuentra en la fase de \textbf{Prueba Cerrada} dentro de Google Play Console, por lo que solo un grupo de evaluadores (\textit{testers}) tienen acceso para probarla. Esto permite observar cómo funciona en diferentes dispositivos móviles y asegurar que toda la tecnología PWA y TWA responde correctamente antes de lanzarla de forma pública para todos los usuarios.